\paragraph{Stationary evolution in the full volume.}
The preceding interpolation need not be stationary. A different, explicitly
declared initial-value problem uses every volume edge as an independent
coordinate. In this paragraph write \(D=\widetilde D\),
\(C=\widetilde C\), \(M=M_1\), \(K=C^{\mathsf T}M_2C\) and
\(F=M+h^2K/6\). These are the geometric matrices above, not counting
pairings. Since \(M>0\), \(K\geq0\) and \(CD=0\), one has
\(F>0\) and \(KD=0\). Sources have the same covector and temporal
impulse interpretation as before.

\begin{theorem}[Full-volume prism evolution and constraint preservation]
\label{thm:whitney-volume-dynamics}
For \(h>0\), the interior vector-potential equations of the exactly
integrated prism action are equivalent to
\begin{equation}
 F A_{n+1}=\left(2M-\frac{2h^2}{3}K\right)A_n-F A_{n-1}
 +h^2J_{n-1}-hMD(\phi_n-\phi_{n-1}).
 \label{eq:whitney-prism-recurrence}
\end{equation}
They determine a unique next vector potential for every supplied scalar
potential and pair of previous vector potentials. With
\(E_n=-(A_{n+1}-A_n)/h-D\phi_n\), the Gauss residual
\(g_n=D^{\mathsf T}ME_n-\rho_n\) satisfies
\[
 g_n-g_{n-1}=-hD^{\mathsf T}J_{n-1}-(\rho_n-\rho_{n-1}).
\]
Thus conserved sources propagate the initial constraint, including any
initial error. For neutral \(\rho_0\), a given initial electric field
\(E_{\rm old}\) can instead be replaced by
\begin{equation}
 E_0=E_{\rm old}+Dz,\qquad
 z=(D^{\mathsf T}MD+\Pi_{13})^{-1}
            (\rho_0-D^{\mathsf T}ME_{\rm old}).
 \label{eq:whitney-gauss-projection}
\end{equation}
This is the minimum \(M\)-norm correction satisfying Gauss's law. It
preserves the curl of the initial electric field and defines a new initial
condition.
\end{theorem}

\begin{proof}
Substitute the definitions of \(E_n\) and \(B_n=CA_n\) in the full
action derivatives above and collect the next-potential coefficient.
Positive definiteness of \(F\) proves existence and uniqueness in the
finite-dimensional edge space. Applying \(D^{\mathsf T}\) to the vector
equation removes every magnetic term and gives the displayed residual
identity. The connected cone has
\(\ker D=\operatorname{span}\{\mathbf1\}\), so
\(D^{\mathsf T}MD\) is invertible on mean-zero vertex vectors.
Neutrality places the right-hand side of
\eqref{eq:whitney-gauss-projection} in that subspace. Any competing
correction differs from \(Dz\) by \(w\) with
\(D^{\mathsf T}Mw=0\); hence \((Dz)^{\mathsf T}Mw=0\).
Pythagoras proves minimality. Finally, \(CDz=0\).
\end{proof}

\begin{theorem}[Work identity and sharp positive-mode stability]
\label{thm:whitney-prism-stability}
Let \(\overline B_n=(B_n+B_{n+1})/2\). Along the full recurrence,
\begin{align*}
 H_n&=\frac12E_n^{\mathsf T}\left(M-\frac{h^2}{12}K\right)E_n
             +\frac12\overline B_n^{\mathsf T}M_2\overline B_n,\\
 H_n-H_{n-1}&=-\frac h2 J_{n-1}^{\mathsf T}(E_n+E_{n-1}).
\end{align*}
For a positive generalized mode \(Kv=\lambda Mv\), all source-free
modal solutions are bounded exactly when \(0<h^2\lambda<12\).
At \(h^2\lambda=12\), the double characteristic root \(-1\) admits
linearly growing alternating solutions. Above twelve an exponentially
growing solution exists. Source-free zero modes have constant electric field and can
have linearly drifting gauge potentials.
\end{theorem}

\begin{proof}
The recurrence also reads
\(F(E_n-E_{n-1})=hKA_n-hJ_{n-1}\); here \(KD=0\)
removes the scalar-potential terms. Faraday's identity is
\(B_{n+1}-B_n=-hCE_n\). Polarize the two quadratic terms, use these
identities, and cancel the magnetic cross terms. Equivalently,
\(H_n=E_n^{\mathsf T}FE_n/2+B_n^{\mathsf T}M_2B_{n+1}/2\),
which makes the telescoping cancellation immediate. This proves the work
law in every gauge.

For \(z=h^2\lambda\), the scalar characteristic equation is
\[
 r^2-2a_zr+1=0,\qquad a_z=\frac{1-z/3}{1+z/6}.
\]
For positive \(z\), the strict inequality \(|a_z|<1\) holds exactly
on \(z<12\). Its two distinct unit-circle roots give bounded solutions.
At twelve, \((-1)^{n+1}n\) is a solution and its electric difference is
unbounded. Above twelve, the negative real reciprocal roots include one
with modulus greater than one. In temporal gauge the zero-mode recurrence
is a second difference, giving affine potentials. These facts do not bound gauge
potentials by field energy.
\end{proof}

The finite action expansions, weak Noether identity, work law, uniqueness
and sharp scalar alternatives are formalized in
\path{Lean/Screen/WhitneyMaxwellDynamics.lean}. Stable initial-value
evolution does not imply a nonsingular fixed-endpoint boundary problem:
for two slabs and one mode with \(h^2\lambda=3\), zero endpoint
potentials leave the interior potential arbitrary when its source is zero,
and are incompatible with a nonzero interior source. The initial-value
coefficient \(1+h^2\lambda/6\) remains positive there.

\paragraph{An exact bound for the supplied cone.}
Every tetrahedron has three apex rays with Gram matrix
\(2I_3+\varphi\mathbf1\mathbf1^{\mathsf T}\).
Divide each local Whitney mass and curl-energy matrix by its positive
tetrahedron volume. Exact arithmetic in \(\mathbb Q(\sqrt5)\) gives an
\(LDL^{\mathsf T}\) factorization of \(24M_T-K_T\) with strictly
positive diagonal. An independent verifier reconstructs the barycentric
gradients from the supplied vertex coordinates, checks the matrix identity
and encloses \(\sqrt5\) between rational bounds to certify each sign.
The twenty signed local-to-global edge maps cover every global edge, so
assembly proves \(K<24M\), without a floating-point eigenvalue bound.
At \(h=1/2\), therefore,
\[
 M-h^2K/12>M/2,\qquad
 E_n^{\mathsf T}ME_n\leq4H_n,\qquad
 B_n^{\mathsf T}M_2B_n\leq16H_n
\]
on a source-free interval. The magnetic estimate also holds at the other
endpoint of each slab: use
\(B_n=\overline B_n+hCE_n/2\) and
\(h^2E_n^{\mathsf T}KE_n\leq6E_n^{\mathsf T}ME_n\).
The factorization is an exact algebraic certificate; the trajectory below
is a separate numerical calculation.

\paragraph{A recorded stationary volume history.}
Retain the inherited initial vector potential and prescribed neutral source
covectors, apply \eqref{eq:whitney-gauss-projection}, and use
\eqref{eq:whitney-prism-recurrence}. The correction changes the initial
electric field by approximately \(0.71279026\) in \(M\)-norm.
Consequently this history is distinct from the earlier counting-action
history and its cone interpolation. Each of three slices occupies thirteen
scalar and forty-two edge registers. For each edge, read both endpoint
pair averages, retain scalar baselines, and restore the registers by
\((b,2r-b)\). The resulting bounded observer-like software patch has
writable ports, readback, records, feedback and a public evidence bundle.
Its \(805\) recorded events include \(252\) feedback cycles. Every
read names its last writer and exact value; the volume advance consumes
the decoded registers and source covectors. The dense metric projection
and evolution rule are supplied software operations, not consequences of
pair averaging or local causal ancestry.

The two endpoint-fixed gauge representatives have field-plus-source action
approximately \(-1.14785515542\). Independent spacetime quadrature and
complex-step differentiation verify all \(42+2\cdot13=68\) free action
derivatives, including radial edges and the apex. Their absolute residuals
are below \(10^{-12}\). The metric evolution uses double precision and
declared \(10^{-9}\) comparison tolerances; register values are exact
rational encodings of the numerical outputs, and feedback restoration is
exact for those encodings. No interval or exact-real trajectory certificate
is asserted.

A separate \(64\)-slab numerical continuation starts from the decoded
temporal-gauge data. Its maximum Gauss, vector-equation and work residuals
are below \(10^{-12}\), as is the source-free energy drift. Its departure
from the fixed radial extension reaches approximately \(0.45218\),
showing that the volume recurrence exercises independent radial degrees
of freedom. Only the first three slices have the serial execution trace.
The remaining continuation is a numerical field calculation.
Prescribed currents are not dynamical charged matter. Stationarity against
all finite-element coordinates also does not imply stationarity against
arbitrary smooth continuum tests. Spatial refinement, physical clock and
metric selection, and laboratory identification require further input.

\paragraph{The Hilbert space of the same radiative modes.}
Consider the source-free neutral sector of this same geometric action,
in continuous time and temporal gauge, with the spatial mesh fixed. The constraint is
\(D^{\mathsf T}M\dot A=0\). A time-independent gauge transformation
removes the constant longitudinal potential, giving the transverse space
\(V_\perp=\ker(D^{\mathsf T}M)\). The cone is contractible:
its explicit curl blocks imply \(\ker C=\operatorname{im}D\),
of dimension twelve. Thus \(V_\perp\) has dimension thirty and
\(K\) is positive definite there. The finite-dimensional spectral theorem
gives a complete real frame \(v_1,\ldots,v_{30}\) with
\[
 v_i^{\mathsf T}Mv_j=\delta_{ij},\qquad
 Kv_i=\omega_i^2Mv_i,\qquad \omega_i>0.
\]
For \(A=\sum_i q_i v_i\), the exact spatially semidiscrete action is
\[
 S_0=\frac12\int\left(\dot A^{\mathsf T}M\dot A-A^{\mathsf T}KA\right)dt
    =\frac12\int\sum_i(\dot q_i^2-\omega_i^2q_i^2)\,dt.
\]
The spatial edge-element action is standard \cite{HeTeixeira2005}, as are
the finite normal-mode theorem and canonical field quantization. The
attachment here is to these geometric matrices.

\begin{theorem}[Canonical quantum realization of the volume radiative sector]
\label{thm:whitney-fock-bridge}
Supply canonical bosonic quantization and \(\hbar>0\). The preceding
sector has Hilbert space \(\mathcal H=\ell^2(\mathbb N^{30})\).
On its finite-occupation domain, let
\begin{align*}
 a_i|n\rangle&=\sqrt{n_i}\,|n-e_i\rangle,&
 a_i^\dagger|n\rangle&=\sqrt{n_i+1}\,|n+e_i\rangle,\\
 \widehat q_i&=\sqrt{\frac{\hbar}{2\omega_i}}(a_i+a_i^\dagger),&
 \widehat p_i&=i\sqrt{\frac{\hbar\omega_i}{2}}(a_i^\dagger-a_i).
\end{align*}
Then \([\widehat q_i,\widehat p_j]=i\hbar\delta_{ij}I\), and
\[
 \widehat H=\frac12\sum_i(\widehat p_i^2+\omega_i^2\widehat q_i^2),
 \qquad E_n=\hbar\sum_i\omega_i(n_i+1/2).
\]
The diagonal Hamiltonian on
\(\mathcal D(\widehat H)=\{\psi:\sum_n E_n^2|\psi_n|^2<\infty\}\)
is self-adjoint and generates the strongly continuous unitary group
\((U_t\psi)_n=e^{-itE_n/\hbar}\psi_n\).
The Heisenberg equations are
\(\dot{\widehat q}_i=\widehat p_i\) and
\(\dot{\widehat p}_i=-\omega_i^2\widehat q_i\).
Consequently \(\widehat A=\sum_i v_i\widehat q_i\),
\(\widehat E=-\sum_i v_i\widehat p_i\) and
\(\widehat B=C\widehat A\) use the same classical mode shapes and
satisfy the source-free semidiscrete Maxwell equations on the common
invariant domain.
\end{theorem}

\begin{proof}
Start with polynomials in thirty complex variables. Give monomials the
inner product \(\langle X^n,X^m\rangle=n!\delta_{nm}\),
where \(n!=\prod_i n_i!\). Multiplication by \(X_i\) and
differentiation by \(X_i\) are adjoint on this domain and satisfy
\([a_i,a_j^\dagger]=\delta_{ij}I\). Normalizing monomials gives the
displayed orthonormal occupation basis and its \(\ell^2\) completion.
The quadratic Hamiltonian and commutators follow by expansion. Finite
occupation vectors are analytic for each position and momentum operator.
If \(N\) bounds a vector's occupation support, its \(k\)-th iterate
has norm at most \(C_\psi C^k\sqrt{(N+k)!/N!}\), with fixed
constants \(C_\psi,C\). This makes the analytic-vector series converge
for sufficiently small argument. Thus these symmetric operators have
unique self-adjoint closures by the analytic-vector theorem
\cite{Nelson1959}.

The adjoint of the real diagonal operator has exactly the stated weighted
square-summability domain, proving self-adjointness of \(\widehat H\).
Truncation in the occupation basis proves that finite-support vectors are
a core. The same finite-tail argument proves strong continuity of
\(U_t\). Finally
\([\widehat H,\widehat q_i]=-i\hbar\widehat p_i\) and
\([\widehat H,\widehat p_i]=i\hbar\omega_i^2\widehat q_i\)
give the equations, and the complete spectral frame pulls them back to
the actual mass and stiffness forms.
\end{proof}

\path{Lean/Screen/WhitneyQuantumBridge.lean} formalizes the complete
normal-frame action pullback, polynomial commutation relations, occupation
energies and Heisenberg identities. Its frame is supplied; existence here
uses the finite spectral argument above. The Hilbert completion and
self-adjoint-domain argument are the analytic proof just given, not Lean
results. This is canonical quantization of the free radiative sector.
The recorded history with nonzero prescribed charge additionally needs an
affine longitudinal sector and source coupling. Classical readback does
not select a quantum state, Born statistics, \(\hbar\) or a physical
clock, and the finite mesh supplies no relativistic quantum field net.

\paragraph{A controlled temporal limit.}
The finite-step prism dynamics and the continuous-time oscillator have
different frequencies. For one mode, put \(z=h^2\lambda\),
\(f=1+z/6\) and choose \(0<z<12\). The discrete Legendre momenta
give the exact symplectic map
\[
 \begin{pmatrix}q'\\p'\end{pmatrix}
 =\frac1f\begin{pmatrix}
 1-z/3&h\\-h\lambda(1-z/12)&1-z/3
 \end{pmatrix}\begin{pmatrix}q\\p\end{pmatrix}.
\]
Define \(\theta\in(0,\pi)\) by
\(\cos\theta=(1-z/3)/f\). Since
\(\sin^2\theta=z(1-z/12)/f^2\), this is the time-\(h\)
flow of the positive modified Hamiltonian
\[
 H_h=\frac{\theta}{2f\sin\theta}
        \left[p^2+\lambda(1-z/12)q^2\right].
\]
On any fixed bounded positive spectral window,
\(\theta/h=\sqrt\lambda+O(h^2)\) and the quadratic coefficients
of \(H_h-H_0\) are \(O(h^2)\), uniformly as \(h\to0\).
On bounded time intervals the corresponding classical propagators
therefore converge at order two for bounded canonical initial data.
The initial pair of potentials must be obtained from the discrete
Legendre map for those data; an arbitrary first-order initial displacement
does not provide the same order claim. This is a finite-dimensional
temporal convergence argument with the spatial mesh fixed.
Canonical quantization of \(H_h\) supplies its own unitary oscillator
flow, with frequency \(\theta/h\); its continuous Hamiltonian path
fixes the metaplectic lift. Neither its generator nor its frequency equals
\(\widehat H\) at fixed \(h\). No uniform operator-norm error bound
for the unbounded quantum generators, spatial convergence theorem or
physical time calibration follows from this calculation.
